%% SPDX-License-Identifier: CC-BY-4.0
%% SPDX-FileCopyrightText: 2026 Serge TABE DJATO
% !TeX program = xelatex
%% ============================================================
%% IntrepidCore — SPÉCIFICATION ICES v0.1
%% Le MissionBundle, la canonicalisation, le sceau Ed25519
%% ============================================================
\documentclass[11pt,a4paper]{article}
\usepackage[a4paper,top=2.4cm,bottom=2.6cm,left=2.5cm,right=2.5cm]{geometry}
\usepackage[french]{babel}\usepackage{csquotes}\usepackage{fontspec}\usepackage{microtype}
\IfFontExistsTF{Times New Roman}{\setmainfont{Times New Roman}}{\setmainfont{TeX Gyre Termes}}
\usepackage{amsmath}\usepackage{graphicx}\usepackage{booktabs}\usepackage{array}
\usepackage{enumitem}\usepackage{titlesec}\usepackage{xcolor}\usepackage{fancyhdr}\usepackage{tcolorbox}
\usepackage{listings}\usepackage{tikz}\usetikzlibrary{arrows.meta,positioning}
\tcbuselibrary{skins}\usepackage{hyperref}

\definecolor{intrepid}{RGB}{0,57,107}\definecolor{accent}{RGB}{213,93,36}
\definecolor{cyanIC}{RGB}{0,196,212}\definecolor{vertAtlas}{RGB}{64,145,108}
\definecolor{rowgray}{RGB}{240,240,240}\definecolor{codebg}{RGB}{245,247,250}
\IfFontExistsTF{Virgil 3 YOFF}{\newfontfamily\chartfont{Virgil 3 YOFF}}{\let\chartfont\sffamily}

\lstset{basicstyle=\ttfamily\footnotesize, backgroundcolor=\color{codebg},
  breaklines=true, frame=single, rulecolor=\color{gray!40}, columns=fullflexible,
  keywordstyle=\color{intrepid}\bfseries, commentstyle=\color{vertAtlas}}

\hypersetup{hidelinks, pdfauthor={IntrepidCore}, pdftitle={ICES v0.1 — Spécification}}

\newcommand{\mylogo}{\includegraphics[height=2cm]{logo}}
\newcommand{\headerblock}{%
  \noindent\begin{minipage}[c]{0.45\linewidth}\mylogo\end{minipage}\hfill
  \begin{minipage}[c]{0.52\linewidth}\raggedleft
    {\small\color{intrepid}\textit{Innover localement. Impacter globalement.}}\\[4pt]
    {\small\color{intrepid}Build intelligent engineering for African infrastructure.}\\[4pt]
    {\footnotesize\color{gray}Spécification ICES v0.1 \textperiodcentered\ Juin~2026 \textperiodcentered\ \ttfamily contact@intrepidcore.io}%
  \end{minipage}\par\vspace{5pt}\rule{\linewidth}{1pt}\vspace{0.4cm}}

\newtcolorbox{icesbox}[2][]{enhanced, colback=#2!8!white, colframe=#2, boxrule=0.6pt,
  left=6pt, right=6pt, top=5pt, bottom=5pt, fonttitle=\bfseries\small\color{#2}, title={#1}}

\setlength{\parindent}{0pt}\setlength{\parskip}{6pt plus 1pt minus 0.5pt}
\titleformat{\section}{\normalfont\fontsize{12.5}{14}\bfseries\color{intrepid}}{\thesection.}{0.6em}{}
\titleformat{\subsection}{\normalfont\fontsize{11}{13}\bfseries\color{accent}}{}{0em}{}

\pagestyle{fancy}\fancyhf{}
\fancyhead[L]{\small\color{intrepid}\textsc{IntrepidCore} \textemdash\ Spécification ICES v0.1}
\fancyhead[R]{\small\thepage}
\fancyfoot[C]{\small\textit{Standard ouvert \textemdash\ CC BY 4.0, \textcopyright{} 2026 Serge TABE DJATO \textemdash\ Diffusion libre}}
\renewcommand{\headrulewidth}{0.4pt}\renewcommand{\footrulewidth}{0.4pt}

\begin{document}
\thispagestyle{empty}
\headerblock

\begin{center}
  {\fontsize{16}{19}\bfseries\color{intrepid} SPÉCIFICATION ICES v0.1}\\[4pt]
  {\fontsize{11}{14}\bfseries Le MissionBundle, la canonicalisation et le sceau Ed25519}\\[4pt]
  {\fontsize{10}{12}\normalfont\color{gray} Standard ouvert \textperiodcentered\ Juin~2026 \textperiodcentered\ Artefacts~: \texttt{mission-bundle.schema.json}, \texttt{ices\_verify.py}}
\end{center}
\vspace{0.2cm}\rule{\linewidth}{0.5pt}\vspace{0.2cm}

\begin{icesbox}[Objet]{intrepid}
ICES (IntrepidCore Engineering Standard) définit un \textbf{format de preuve d'intégrité de
calcul d'ingénierie}. Un livrable ICES est un \texttt{MissionBundle} scellé par signature
\textbf{asymétrique Ed25519}, qu'un tiers (bailleur, auditeur) vérifie \textbf{hors-ligne, sans
secret partagé}. ICES est un \textbf{standard ouvert et gratuit}~; sa génération et sa
vérification ne donnent lieu à aucune facturation.
\end{icesbox}

\section{Le MissionBundle}
Objet JSON typé (schéma normatif~: \texttt{mission-bundle.schema.json}, JSON~Schema 2020-12),
en cinq sections~: \texttt{header}, \texttt{inputs}, \texttt{scenario} (optionnel),
\texttt{outputs}, \texttt{gates}, plus le sceau \texttt{ices\_envelope}. Un bundle \emph{sans}
\texttt{scenario} ni \texttt{ices\_envelope} est l'état legacy déterministe~:
\textbf{rétrocompatibilité par construction}.

\begin{icesbox}[Enveloppe universelle vs charge utile]{accent}
\texttt{header}, \texttt{gates} et \texttt{ices\_envelope} sont \textbf{identiques quel que soit
le produit émetteur} (LCPI, Atlas, IgnisCore). \texttt{inputs} et \texttt{outputs} portent une
partie universelle (config, FallbackRegistry) et une partie \textbf{propre au domaine}, validée
par des règles conditionnelles (\texttt{if header.domain == \ldots then \ldots}) dans le schéma
normatif. C'est ce qui rend ICES nativement extensible à Atlas et IgnisCore sans dénaturer sa
structure hydraulique d'origine.
\end{icesbox}

\begin{center}\small
\begin{tabular}{@{}>{\ttfamily}p{2.6cm} p{11.6cm}@{}}
\toprule
\normalfont\textbf{Section} & \textbf{Contenu (résumé)} \\
\midrule
header & schema\_version, project\_key, \textbf{domain} (aep/hydro/atlas/igniscore/\ldots), \textbf{engine\_version} (identifiant neutre du moteur émetteur), date \\
inputs & \textbf{universel}~: config canonisée, FallbackRegistry \textperiodcentered\ \textbf{spécifique au domaine}~: réseau (aep/hydro), geo\_refs (IDs réalisations Atlas), geostatistics (atlas), biomass\_mix (igniscore) \\
scenario & ensemble\_seed, n\_realizations, axes d'incertitude (sol, demande, IWS, coupures) — non applicable à \texttt{atlas} (déjà probabiliste via SGS) \\
outputs & \textbf{spécifique au domaine}~: distribution/coûts (aep/hydro), realizations (atlas), biocore (igniscore) \textperiodcentered\ robustness (P(service), équité) quand l'Ensemble Engine s'applique \\
gates & les contrôles EPS-001 (PASS/WARN/FAIL) — la preuve inclut aussi les échecs \\
ices\_envelope & alg=Ed25519, digest SHA3-256, sig, \textbf{pubkey} (clé publique vérifiable) \\
\bottomrule
\end{tabular}
\end{center}

\section{Canonicalisation déterministe}
\textit{Règle d'or~: même entrée logique $\rightarrow$ mêmes octets, sur toute machine.} Sinon le
hash dérive et la preuve s'effondre.
\begin{itemize}[leftmargin=1.3em, itemsep=1pt]
  \item sérialisation JCS (RFC~8785)~: clés d'objet \textbf{triées}~; séparateurs compacts ; UTF-8 ;
  \item flottants \textbf{arrondis} à une précision fixée par champ ;
  \item chemins machine $\rightarrow$ \textbf{empreintes} (\texttt{sha}), jamais les chemins absolus ;
  \item géométrie \textbf{snappée} au préalable (logique de \texttt{canonize.py} généralisée).
\end{itemize}

\section{Le sceau}
Au point de sortie unique du moteur (\texttt{finalize()}), quel que soit le produit~:
\[
\texttt{digest} = \text{SHA3-256}\big(\,\text{canon(inputs)}\ \|\ \text{canon(scenario)}\ \|\
\text{engine\_version}\ \|\ \text{canon(outputs)}\ \|\ \text{canon(gates)}\,\big)
\]
\[
\texttt{sig} = \text{Ed25519.sign}(\texttt{digest},\ sk_{\text{émetteur}})
\]
La clé privée $sk$ reste chez l'émetteur~; la clé publique $pk$ est diffusée. \textbf{Ed25519}
(RFC~8032) est choisi pour son \textbf{déterminisme} (sceau reproductible) et son indépendance à
tout aléa de signature.

\begin{figure}[h]\centering
\begin{tikzpicture}[font=\chartfont\footnotesize, node distance=0.9cm,
  b/.style={draw=intrepid, fill=intrepid!8, rounded corners=2pt, inner sep=4pt, text=intrepid}]
\node[b] (i) {inputs\,$\|$\,scenario\,$\|$\,engine\,$\|$\,outputs\,$\|$\,gates};
\node[b, right=1.1cm of i, draw=accent, text=accent, fill=accent!10] (h) {SHA3-256};
\node[b, right=1.0cm of h, draw=cyanIC, text=black, fill=cyanIC!12] (s) {Ed25519 sign};
\node[b, right=1.0cm of s, fill=intrepid!14] (e) {ices\_envelope};
\node[draw=vertAtlas, fill=vertAtlas!12, rounded corners=2pt, below=0.7cm of s, inner sep=4pt,
   text width=5cm, align=center] (v) {\textbf{\color{vertAtlas}ices verify} — clé publique, hors-ligne};
\draw[-{Stealth[length=2mm]}, intrepid] (i)--(h); \draw[-{Stealth[length=2mm]}, intrepid] (h)--(s);
\draw[-{Stealth[length=2mm]}, intrepid] (s)--(e);
\draw[-{Stealth[length=2mm]}, vertAtlas] (e.south) |- (v.east);
\end{tikzpicture}
\end{figure}

\section{Vérification (\texttt{ices verify})}
Le vérificateur de référence (\texttt{ices\_verify.py}~; cible production~: crate Rust
\texttt{ices-core}, binaire offline) exécute trois étapes~: (1) recanonicalise les sections
signées~; (2) recalcule le condensat SHA3-256~; (3) vérifie la signature Ed25519 avec la clé
publique. Toute modification d'un extrant après signature fait diverger le condensat
$\rightarrow$ \textbf{INVALIDE}. Le livre blanc ICES v4.0 ajoute un contrôle~: la clé publique doit figurer
dans une liste de clés de confiance, valide et non révoquée à la date de signature. Le vérificateur de
référence l'implémente et refuse par défaut un sceau si aucune liste ne lui est fournie.

\begin{icesbox}[Preuve de concept (banc de test)]{vertAtlas}
Banc exécuté~: un \texttt{MissionBundle} signé est déclaré \texttt{VALIDE}~; le même bundle dont
on baisse le CAPEX \emph{après} signature est déclaré \texttt{INVALIDE} (intégrité~: ÉCHEC). La
fraude post-calcul est détectée par un tiers, sans secret. \textit{C'est la première des quatre
preuves de structuration (jalon~G1).}
\end{icesbox}

\section{Extension multi-domaine (Atlas, IgnisCore)}
Le schéma v2 est délibérément \textbf{généralisé}~: seuls \texttt{header}, \texttt{gates} et
\texttt{ices\_envelope} sont universels~; \texttt{inputs}/\texttt{outputs} varient par domaine
via des règles conditionnelles (\texttt{allOf/if-then} sur \texttt{header.domain}). Un bundle
\texttt{atlas} exige \texttt{inputs.geostatistics} et \texttt{outputs.realizations}~; un bundle
\texttt{igniscore} exige \texttt{inputs.biomass\_mix} et \texttt{outputs.biocore}~; un bundle
\texttt{aep/hydro} exige le réseau et la distribution.

\begin{icesbox}[Preuve de généralisation (double banc de test)]{cyanIC}
Un \texttt{MissionBundle} de domaine \texttt{atlas} (run géostatistique~: variogramme, méthode
SGS, IDs de sondage en entrée~; réalisations P10--P90 en sortie) a été signé, validé contre le
schéma et vérifié \texttt{VALIDE}, au même titre qu'un bundle \texttt{aep}. Contre-épreuve~:
maquiller ce même bundle en \texttt{domain: aep} le fait \textbf{rejeter} par le schéma (champs
hydrauliques requis absents) — la séparation enveloppe/charge utile fonctionne, elle n'est pas
qu'une déclaration d'intention.
\end{icesbox}

\begin{icesbox}[Limite assumée — ce qui n'est pas encore fait]{accent}
Cette généralisation porte le \textbf{schéma et le vérificateur}. Les codebases Atlas et
IgnisCore n'exposent pas encore de point \texttt{finalize()} équivalent à celui de LCPI~: le
câblage réel (émission d'un bundle scellé à la fin d'un run de krigeage ou d'un dimensionnement
BioCore) reste un chantier d'intégration à part entière, non encore réalisé. Ne pas
surpromettre~: le standard est prêt à les recevoir, il ne les couvre pas encore en production.
\end{icesbox}

\section{Statut \& gouvernance}
ICES v0.1 est publié en \textbf{standard ouvert}. La spécification, le schéma JSON et le
vérificateur de référence sont libres d'accès. Gestion des clés~: HSM/KMS pour la racine, clés
FDE par mission. ICES ne certifie pas la responsabilité de l'ingénieur, mais l'\textbf{intégrité}
du calcul (intrants, moteur, extrants).

\vspace{0.3cm}\rule{\linewidth}{0.4pt}
\begin{center}\footnotesize\color{gray}
ICES v0.1 \textperiodcentered\ IntrepidCore \textperiodcentered\ Artefacts normatifs~:
\texttt{mission-bundle.schema.json} \textperiodcentered\ \texttt{ices\_verify.py}
\end{center}
\end{document}
